\documentclass[18pt]{article}

\usepackage[margin=.75in]{geometry}


\title{A Bayesian Framework for Bathymetry Estimation from Surface Measurements of a Shallow Fluid}
\author{Patricio Clark Di Leoni, Nathan Glatt-Holtz\\
  \scriptsize{emails: negh@iu.edu, pclarkdileoni@udesa.edu.ar}}

%\date{}

\usepackage{amsfonts, amssymb, amsmath, amsthm, multicol}
\usepackage[colorlinks=true, pdfstartview=FitV, linkcolor=blue,
            citecolor=blue, urlcolor=blue, bookmarksdepth=3]{hyperref}
\usepackage[usenames]{color}
\definecolor{Red}{rgb}{0.7,0,0.1}
\definecolor{Green}{rgb}{0,0.7,0}
\usepackage{accents}
\usepackage{comment}
\usepackage{float}
\usepackage{graphicx}
\usepackage{multirow}
\usepackage{parcolumns}
\usepackage{subcaption}
\usepackage{textcomp}
\usepackage{mathrsfs}
\usepackage{mathtools}
\usepackage{dsfont}
\usepackage{nccmath}
\usepackage{enumerate}
\usepackage[]{mdframed} \usepackage{cjhebrew}
\usepackage[capitalize,nameinlink,noabbrev]{cleveref}

\providecommand*\showkeyslabelformat[1]{{\normalfont \tiny#1}}
\usepackage[notref,notcite,color]{showkeys}
\definecolor{labelkey}{rgb}{0,0,1}

\pagestyle{myheadings}
\numberwithin{equation}{subsection}

%%%% MACROS %%%%%%%%%%%%%%%%%%%%%%%%%%%%%%%%%%

\newtheorem{Thm}{Theorem}[section]
\newtheorem{Lem}[Thm]{Lemma}
\newtheorem{Prop}[Thm]{Proposition}
\newtheorem{Cor}[Thm]{Corollary}
\newtheorem{Ex}[Thm]{Example}
\newtheorem{Def}[Thm]{Definition}
\newtheorem{Rmk}[Thm]{Remark}
\newtheorem{Conj}[Thm]{Conjecture}
\newtheorem{Hypo}[Thm]{Hypothesis}

\newtheorem{Not}{Notation}[section]
\newtheorem*{LemNoNum}{Lemma}
\newtheorem*{Thm*}{Theorem}

%Internal Editing
\newcommand{\alert}[1]{\textcolor{red}{#1}}

\newcommand{\RR}{\mathbb{R}}
\newcommand{\indFn}{1 \! \! 1}
\newcommand{\Prb}{\mathbb{P}}
\newcommand{\EE}{\mathbb{E}}

\newcommand{\beps}{\boldsymbol{\epsilon}}
\newcommand{\Obs}{\mathcal{O}}
\newcommand{\bfY}{\mathbf{Y}}



%%%%%%

\begin{document}
  \markboth{(DRAFT COPY-- Please do not distribute)}
{(DRAFT COPY-- Please do not distribute)}

\maketitle

\begin{abstract}
\end{abstract}

{\noindent \small {\it \bf Keywords:} \\
  {\it \bf MSC2020:}}

\newpage

%\begin{footnotesize}
\setcounter{tocdepth}{2}
\tableofcontents
%\end{footnotesize}




\newpage

%---------------------------------------------------------
% Introduction
%---------------------------------------------------------

\section{Introduction}\label{sec:intro}


We are interested in the estimation of the Bathymetry of shallow water flow from error corrupted surface wave height measurements. 
The equations of motion are as follows:
\begin{align}
    &h_t + (hu)_x + (hv)_y = 0
    \label{eq:sw:h}\\
   &(hu)_t + \bigl(hu^2 + \tfrac{1}{2}gh^2\bigr)_{\!x} + (huv)_y = -ghb_x
    \label{eq:sw:u}\\
   &(hv)_t + (huv)_x + \bigl(hv^2 + \tfrac{1}{2}gh^2\bigr)_{\!y} = -ghb_y
   \label{eq:sw:v}
\end{align}
where $h$ is the water column height, $(u,v)$ is the horizontal velocity and $g$ is the 
gravitational constant.  We supplement \eqref{eq:sw:h}--\eqref{eq:sw:v} with
initial conditions 
\begin{align}
	h(0,x,y) = h_0(x,y), \quad u(0,x,y) = u_0(x,y), \quad v(0,x,y) = v_0(x,y),
\end{align}
and boundary conditions.\footnote{The fluid domain $D$ of definition needs to be defined and in general $D$ is not fixed. Even if
we imposed e.g. periodic boundary conditions the fluid motion here is hyperbolic and can form
shocks in finite time.   At least from an analysis point of view we may wish to 
impose further conditions (small data? viscous dissipation?) to obtain a more mathematically (and numerically?) tractable set 
of equations.}   

We will measure the surface elevation $\eta$ that is
\begin{align}
	\eta(t,x,y) :=  h(x,y,t)+b(x,y).
\end{align}
sparsely and with corrupting noise in our observations.  More specifically our model for a set of observation 
$\bfY= (y_1, …, y_n)$ is 
as follows
\begin{align}
\label{eq:obs:proceedure}
y_i = \eta(t_i, x_i, y_i; b) + \epsilon_i
\end{align}
where $((t_1, x_1, y_1),\ldots, (t_n, x_n, y_n))$ are the observation locations and $\beps = (\epsilon_1, \ldots, \epsilon_n)$ is an (additive) error or noise
corruption in the measurement of $(\eta(t_1, x_1, y_n; b), \ldots,  \eta(t_n, x_n, y_n; b))$.\footnote{More generally we can just as well consider $y_i = \Obs_j(\eta(b)) + \epsilon_i$ where the $\Obs_j$'s are observation
operators.  These $\Obs_j$ could be spatial temporal averages of $\eta$.  Of course the special case of spatial-temporal point observations is given by $\Obs_j(\eta) = \eta(t_j, x_j, y_j)$.}
Suppose that the $\beps \sim N(0, \Sigma)$, this includes the possibility of correlations in the observation noise, we have a the likelihood of $b$


We would be treating the problem of resolving the Bayesian posterior for the bottom boundary given data points $\bfY= (y_1, …, y_n)$ as 
\begin{align}
\mu(d b) \propto \exp\left( -\frac{1}{2} | \Sigma^{1/2} (Y - H(b)) |^2)\right) \mu_0(db),
\end{align}
where 
\begin{align}
	H(b) = (\eta(t_1, x_1, y_n; b), \ldots,  \eta(t_n, x_n, y_n; b)).
\end{align}
Here we are supposing that 



where $\beps =  (\epsilon_1,\ldots, \epsilon_n) \sim N(0,\Sigma), h(\cdot, h_b)$ is the solution for the height of shallow water with
$H(h_b) = (h(t_1, x_1, h_b) , … , h(t_1, x_1, h_b))$.  $\mu_0$ is the prior which we will take to be a Gaussian field with mean $m_b$
and covariance $C$.    

The modeling questions are:
1) what we should take $\epsilon$ the observational uncertainties? One might expect correlations between the measurement error for nearby measurement in space time but perhaps I’m being too fancy galaxy brained here.
2) What is the the mean M and covariance structure of the prior.


\newpage
\section*{Acknowledgements}

Our efforts are supported under the grants  DMS-2108790
(NEGH). 

%\begin{footnotesize}
%\addtocontents{toc}{\protect\setcounter{tocdepth}{1}}
%\addcontentsline{toc}{section}{References}
\bibliographystyle{alpha}
\bibliography{otrefs}
%\end{footnotesize}


\newpage









\begin{multicols}{2}



  
\noindent
Nathan E. Glatt-Holtz\\ {\footnotesize
Department of Statistics and Department of Mathematics\\
Indiana University, Bloomington\\
Web: \url{https://negh.pages.iu.edu}\\
Email: \url{negh@iu.edu}} \\[.2cm]
\columnbreak

\noindent
Andrew Warren\\
{\footnotesize
University of British Columbia\
Department of Mathematics\\
Web: \url{https://andrew-warren.github.io}\\
Email: \url{awarren@math.ubc.ca}}
\end{multicols}









\end{document}
